%% kapa: exam=1 round=33 year=2022 period=1 form=A subject=economics subject_ko=경제학원론 seq=32 number=72
%% source: Q-Net 감정평가사 기출문제 (cst00309) artlSeq=5209997 file=raw/1008/5209997/2022년도 제33회 감정평가사 1차시험 1교시(민법,경제학,부동산학).pdf
%% conversion: kordoc 4.24.51 -> tools/gen1.py (auto); review=proofread-sonnet; flags=none
\begin{question}{72}
경기변동이론에 관한 설명으로 옳은 것은?
\begin{choices}
  \item 신케인즈 학파(new Keynesian)는 완전경쟁적 시장구조를 가정한다.
  \item 신케인즈 학파는 총수요 외부효과(aggregate-demand externality)를 통해 가격경직성을 설명한다.
  \item 신케인즈 학파는 총공급 충격이 경기변동의 근본 원인이라고 주장한다.
  \item 실물경기변동이론은 실질임금의 경직성을 가정한다.
  \item 실물경기변동이론에 따르면 불경기에는 비용 최소화가 달성되지 않는다.
\end{choices}
\end{question}
\answer{2}
